\section{Theory and Analytical Design}
\label{sec:theory-and-design}

The initial geometry follows the conventional transmission-line approximation for the
dominant rectangular-patch resonance \cite{stutzman2013antenna}. The free-space
wavelength at the design frequency is
\begin{equation}
\lambda_0=\frac{c}{f_0}=327.64~\text{mm}.
\label{eq:lambda0}
\end{equation}

\subsection{Patch Width}

The initial width is
\begin{equation}
W=\frac{c}{2f_0}\sqrt{\frac{2}{\varepsilon_r+1}},
\label{eq:patch-width}
\end{equation}
which gives
\begin{equation}
W=129.51~\text{mm}.
\end{equation}

\subsection{Effective Permittivity and Fringing Extension}

The effective dielectric constant is estimated as
\begin{equation}
\varepsilon_{\mathrm{eff}}
=\frac{\varepsilon_r+1}{2}
+\frac{\varepsilon_r-1}{2}
\left(1+12\frac{h}{W}\right)^{-1/2},
\label{eq:effective-permittivity}
\end{equation}
which yields $\varepsilon_{\mathrm{eff}}=2.1605$.

The fringing-field extension is
\begin{equation}
\Delta L
=0.412h\frac{(\varepsilon_{\mathrm{eff}}+0.3)(W/h+0.264)}
{(\varepsilon_{\mathrm{eff}}-0.258)(W/h+0.8)},
\label{eq:length-extension}
\end{equation}
producing $\Delta L=0.834$~mm.

\subsection{Effective and Physical Length}

The effective resonant length is
\begin{equation}
L_{\mathrm{eff}}=\frac{c}{2f_0\sqrt{\varepsilon_{\mathrm{eff}}}}
=111.45~\text{mm},
\label{eq:effective-length}
\end{equation}
and the physical starting length is
\begin{equation}
L=L_{\mathrm{eff}}-2\Delta L=109.79~\text{mm}.
\label{eq:physical-length}
\end{equation}

\subsection{Feed Width and Inset Depth}

The documented design uses a 4.89~mm microstrip feed width as the nominal
50-$\Omega$ value for the selected substrate. The inset depth is estimated from
\begin{equation}
R_{\mathrm{in}}(y_0)=R_{\mathrm{edge}}
\cos^2\left(\frac{\pi y_0}{L}\right).
\label{eq:inset-resistance}
\end{equation}
Using $R_{\mathrm{edge}}=250~\Omega$, a target resistance of 50~$\Omega$, and the
final patch length $L=108.2$~mm gives
\begin{equation}
y_0=\frac{L}{\pi}\cos^{-1}\sqrt{\frac{50}{250}}
=38.13~\text{mm}.
\label{eq:inset-depth}
\end{equation}
The final CST inset depth is 38.2~mm.

\begin{table}[H]
\centering
\caption{Analytical starting values and final CST dimensions.}
\label{tab:analytical-final-dimensions}
\begin{tabular}{@{}llll@{}}
\toprule
\textbf{Quantity} & \textbf{Symbol} & \textbf{Analytical value} & \textbf{Final CST value} \\
\midrule
Patch width & $W$ & 129.51~mm & 129.5~mm \\
Patch length & $L$ & 109.79~mm & 108.2~mm \\
Effective dielectric constant & $\varepsilon_{\mathrm{eff}}$ & 2.1605 & -- \\
Fringing extension & $\Delta L$ & 0.834~mm & -- \\
Feed width & $W_f$ & Documented 4.89~mm & 4.89~mm \\
Inset depth & $d_{\mathrm{inset}}$ & 38.13~mm & 38.2~mm \\
Inset-notch width & $W_{\mathrm{notch}}$ & 6.89~mm & 6.89~mm \\
\bottomrule
\end{tabular}
\end{table}
